\section{Experiments}
\label{sec:experiments}

We evaluate \mmr against established QD algorithms across multiple dimensions, ablation studies, and analysis of key design choices.

\subsection{Experimental Setup}
\label{sec:setup}

\paragraph{Benchmark task.}
We use the $N$-DOF planar arm~\cite{cully2015robots}, a standard QD benchmark that naturally scales to high-dimensional behavior spaces.
A planar arm with $N$ revolute joints has genome $\theta \in [-\pi, \pi]^N$ (joint angles).
Fitness is negative distance from the end-effector to a target position.
The behavior descriptor is the vector of cumulative joint angles, yielding a $d = N$ dimensional behavior space in $[0, 1]^d$.
We test $N \in \{5, 10, 20, 50, 100\}$ to evaluate dimensionality scaling.

\paragraph{Baselines.}
\begin{itemize}
    \item \map~\cite{mouret2015illuminating}: Grid-based with $m = 3$ bins per dimension. Archive size grows as $3^d$.
    \item \cvt~\cite{vassiliades2018using}: Voronoi tessellation with $K = 1000$ centroids (matching \mmr archive size).
    \item \textbf{Random Search}: Top-$K$ by fitness from all evaluated solutions (no diversity pressure).
\end{itemize}

\paragraph{Metrics.}
\begin{itemize}
    \item \textbf{QD-Score@$K$}: Sum of fitness values for the top-$K$ solutions, enabling fair comparison across algorithms with different archive sizes.
    \item \textbf{Uniformity CV}: Coefficient of variation of $k$-nearest-neighbor distances ($k = 5$). Lower values indicate more uniform coverage. $\text{CV} = 0$ is perfectly uniform.
    \item \textbf{Mean fitness}: Average fitness across archive members.
    \item \textbf{Mean pairwise distance}: Average pairwise distance between archive members.
\end{itemize}

\paragraph{Hyperparameters.}
Unless otherwise noted: archive size $K = 1000$, generations $G = 2000$, batch size $B = 200$, $\lambda = 0.5$, mutation $\sigma_{\text{mut}} = 0.1$.
All experiments use 10 independent seeds.
We report mean $\pm$ standard deviation and 95\% confidence intervals.

\subsection{Dimensionality Scaling}
\label{sec:scaling}

Figure~\ref{fig:scaling} shows performance as behavior space dimension increases from 5 to 100.

\begin{figure}[t]
    \centering
    \includegraphics[width=\textwidth]{figures/dimensionality_scaling.pdf}
    \caption{Performance vs.\ behavior space dimension. \textbf{Left}: QD-Score@$K$. \textbf{Center}: Uniformity CV (lower is better). \textbf{Right}: Mean fitness. \mmr maintains strong uniformity as dimension increases, while \map degrades rapidly. Error bars show 95\% CI over 10 seeds.}
    \label{fig:scaling}
\end{figure}

At $d = 20$, \mmr achieves a uniformity CV of $0.066$ compared to $0.832$ for \map ($12\times$ improvement) and $0.417$ for \cvt ($6\times$ improvement).
\map's grid becomes exponentially sparse: $3^{20} \approx 3.5 \times 10^9$ cells for $K = 1000$ solutions.
\cvt avoids this explosion but its centroid-based partitioning still struggles with high-dimensional geometry.
\mmr, operating purely on pairwise distances, is unaffected by dimension.

At $d = 5$, all algorithms perform similarly on QD-Score, though \mmr shows slightly worse uniformity than Random Search.
This is expected: in low dimensions with $K = 1000$ solutions, even random points are well-spread, and \mmr's fitness term ($\lambda = 0.5$) causes mild clustering around fitness optima (Section~\ref{sec:discussion}).

\subsection{Lambda Ablation}
\label{sec:lambda}

Figure~\ref{fig:lambda} shows the effect of $\lambda$ on the fitness--diversity tradeoff at $d = 20$.

\begin{figure}[t]
    \centering
    \includegraphics[width=0.7\textwidth]{figures/lambda_ablation.pdf}
    \caption{Effect of $\lambda$ on the fitness--diversity tradeoff ($d = 20$). Each point is a $\lambda$ value; $x$-axis is mean fitness, $y$-axis is uniformity CV (lower is better). $\lambda = 0.5$ provides a good balance. 5 seeds per configuration.}
    \label{fig:lambda}
\end{figure}

At $\lambda = 0$ (pure fitness), the archive collapses to a cluster of high-fitness solutions with poor diversity (high CV).
At $\lambda = 1$ (pure diversity), solutions spread maximally but fitness degrades.
$\lambda = 0.5$ (default) achieves a favorable balance, and performance is robust across $\lambda \in [0.3, 0.7]$.

\subsection{Archive Size Sensitivity}
\label{sec:archsize}

Table~\ref{tab:archsize} shows performance across archive sizes $K \in \{100, 500, 1000, 2000, 5000\}$ at $d = 20$.

\begin{table}[t]
\centering
\caption{Archive size sensitivity for \mmr at $d = 20$. Mean $\pm$ std over 5 seeds.}
\label{tab:archsize}
\begin{tabular}{r c c c}
\toprule
$K$ & QD-Score@$K$ & Uniformity CV $\downarrow$ & Mean Fitness \\
\midrule
100   & --- & --- & --- \\
500   & --- & --- & --- \\
1000  & --- & --- & --- \\
2000  & --- & --- & --- \\
5000  & --- & --- & --- \\
\bottomrule
\end{tabular}
\end{table}

% Results to be filled from experiments/archive_size_ablation.py

\subsection{Distance Function Comparison}
\label{sec:distance_exp}

We compare Euclidean distance, normalized Euclidean, Gaussian saturating, and exponential saturating distance functions across dimensions $d \in \{5, 10, 20, 50\}$.

\begin{figure}[t]
    \centering
    \includegraphics[width=\textwidth]{figures/distance_comparison.pdf}
    \caption{Distance function comparison across dimensions. Saturating distances improve interior coverage and uniformity, particularly at lower dimensions where Euclidean distance pushes solutions to boundaries.}
    \label{fig:distance}
\end{figure}

% Results to be filled from experiments/distance_comparison.py

\subsection{Runtime Analysis}
\label{sec:runtime}

Table~\ref{tab:runtime} reports wall-clock time per generation for the Rust-based \mmr selection.

\begin{table}[t]
\centering
\caption{Selection time per generation (ms) for \mmr with $K = 1000$, $B = 200$.}
\label{tab:runtime}
\begin{tabular}{r c c}
\toprule
Dimension & Naive (Python) & Lazy Greedy (Rust) \\
\midrule
5   & --- & --- \\
20  & --- & ${\sim}50$ \\
50  & --- & --- \\
100 & --- & --- \\
\bottomrule
\end{tabular}
\end{table}

The lazy greedy implementation achieves $10$--$50\times$ speedup over naive Python selection, with the advantage growing at higher dimensions where the priority queue avoids more redundant distance computations.
